\section{Method}
\label{sec:method}

We consider an evolving knowledge graph observed as a sequence of snapshots $\{G_1, G_2, \ldots, G_T\}$, where each snapshot $G_t = (\mathcal{V}_t, \mathcal{E}_t, \mathcal{R}_t, \mathbf{X}_t)$ consists of entities $\mathcal{V}_t$, typed edges $\mathcal{E}_t \subseteq \mathcal{V}_t \times \mathcal{R}_t \times \mathcal{V}_t$, relation types $\mathcal{R}_t$, and multimodal entity attributes $\mathbf{X}_t$.
Each entity $v \in \mathcal{V}_t$ may carry up to three types of features: structural context from its local neighborhood, textual descriptions $x_v^{\text{text}}$ (e.g., disease definitions, gene summaries), and molecular representations $x_v^{\text{mol}}$ (e.g., SMILES strings converted to Morgan fingerprints).
The benchmark defines a task sequence $\mathcal{T} = \{T_1, T_2, \ldots, T_K\}$ derived from the temporal graph snapshots, where each task $T_k$ provides training triples $(h, r, t) \in \mathcal{E}_k$ for a specific relation type.
The objective is to learn entity embeddings $\mathbf{h}_v \in \mathbb{R}^D$ that perform well on link prediction across all tasks seen so far, without catastrophic forgetting of previously learned relations.

\subsection{Architecture Overview}
\label{subsec:overview}

\ours consists of four stages: (i)~three modality-specific encoders that extract structural, textual, and molecular representations; (ii)~a cross-modal attention fusion module that combines these representations through bidirectional attention; (iii)~a decoder (DistMult) that scores candidate triples using the fused entity embeddings; and (iv)~continual learning mechanisms---modality-aware EWC and multimodal memory replay---that protect previously learned knowledge as the model trains on new tasks.

\subsection{Modality-Specific Encoders}
\label{subsec:encoders}

\paragraph{Structural Encoder.}
We employ a Relational Graph Convolutional Network (R-GCN)~\citep{schlichtkrull2018modeling} with basis decomposition to encode the graph structure.
The message-passing update for entity $v$ at layer $l$ is:
\begin{equation}
    \mathbf{h}_v^{(l+1)} = \sigma \left( \sum_{r \in \mathcal{R}} \sum_{u \in \mathcal{N}_v^r} \frac{1}{|\mathcal{N}_v^r|} \mathbf{W}_r^{(l)} \mathbf{h}_u^{(l)} + \mathbf{W}_0^{(l)} \mathbf{h}_v^{(l)} \right),
    \label{eq:rgcn}
\end{equation}
where $\mathcal{N}_v^r$ denotes the set of neighbors of $v$ under relation $r$, $\mathbf{W}_r^{(l)}$ is the relation-specific weight matrix decomposed as $\mathbf{W}_r^{(l)} = \sum_{b=1}^{B} a_{rb}^{(l)} \mathbf{V}_b^{(l)}$ with $B$ basis matrices, and $\sigma$ is ReLU activation.
We use two R-GCN layers, yielding structural embeddings $\mathbf{h}_v^s \in \mathbb{R}^D$.

\paragraph{Textual Encoder.}
For entities with textual descriptions (disease definitions, gene summaries, biological process annotations), we use a frozen BiomedBERT~\citep{gu2021domain} to extract contextualized token representations, followed by mean pooling and a learned linear projection:
\begin{equation}
    \mathbf{h}_v^t = \mathbf{W}_{\text{proj}} \cdot \text{MeanPool}(\text{BiomedBERT}(x_v^{\text{text}})) + \mathbf{b}_{\text{proj}},
    \label{eq:text}
\end{equation}
where $\mathbf{W}_{\text{proj}} \in \mathbb{R}^{D \times 768}$ projects from BiomedBERT's hidden dimension to the shared embedding dimension $D$.
BiomedBERT weights are frozen to reduce computational cost and avoid catastrophic forgetting of pretrained biomedical knowledge; only the projection layer is trained.
Entities without textual features receive a learned default embedding.

\paragraph{Molecular Encoder.}
Drug nodes carry molecular structure information encoded as Morgan fingerprints~\citep{rogers2010extended} (1024-bit, radius 2), which capture circular substructure patterns.
We pass these through a two-layer MLP with ReLU activation and dropout:
\begin{equation}
    \mathbf{h}_v^m = \text{MLP}(x_v^{\text{mol}}) = \mathbf{W}_2 \cdot \text{ReLU}(\mathbf{W}_1 \cdot x_v^{\text{mol}} + \mathbf{b}_1) + \mathbf{b}_2,
    \label{eq:mol}
\end{equation}
where $\mathbf{W}_1 \in \mathbb{R}^{D \times 1024}$ and $\mathbf{W}_2 \in \mathbb{R}^{D \times D}$.
Entities without molecular features (e.g., diseases, pathways) receive a learned default embedding.

\subsection{Cross-Modal Attention Fusion}
\label{subsec:fusion}

To enable information exchange across modalities, we employ bidirectional cross-modal attention.
Given the three modality embeddings $\mathbf{h}^s, \mathbf{h}^t, \mathbf{h}^m \in \mathbb{R}^{N \times D}$ (where $N$ is the number of entities in the current batch), we compute pairwise cross-attention between all modality pairs.
For modalities $i$ and $j$, the cross-attention output is:
\begin{equation}
    \text{CrossAttn}(\mathbf{h}^i, \mathbf{h}^j) = \text{softmax}\left(\frac{\mathbf{Q}_i \mathbf{K}_j^\top}{\sqrt{D/H}}\right) \mathbf{V}_j,
    \label{eq:crossattn}
\end{equation}
where $\mathbf{Q}_i = \mathbf{h}^i \mathbf{W}_Q^i$, $\mathbf{K}_j = \mathbf{h}^j \mathbf{W}_K^j$, $\mathbf{V}_j = \mathbf{h}^j \mathbf{W}_V^j$ are query, key, and value projections with $H{=}4$ attention heads.
The fused representation combines all cross-attention outputs with residual connections:
\begin{equation}
    \mathbf{h}_v = \text{LayerNorm}\left(\text{MLP}_{\text{fuse}}\left([\mathbf{h}_v^s; \hat{\mathbf{h}}_v^{s \leftarrow t}; \hat{\mathbf{h}}_v^{s \leftarrow m}; \hat{\mathbf{h}}_v^{t \leftarrow s}; \hat{\mathbf{h}}_v^{m \leftarrow s}]\right) + \mathbf{h}_v^s\right),
    \label{eq:fusion}
\end{equation}
where $[\cdot;\cdot]$ denotes concatenation, $\hat{\mathbf{h}}_v^{i \leftarrow j}$ is the cross-attention output from modality $j$ to modality $i$, and MLP$_{\text{fuse}}$ projects the concatenated vector back to dimension $D$.
The residual connection from $\mathbf{h}_v^s$ ensures that the structural signal---available for all entities---serves as the primary representation, with textual and molecular signals providing complementary refinement.

\subsection{Modality-Aware Elastic Weight Consolidation}
\label{subsec:maewc}

Standard EWC~\citep{kirkpatrick2017overcoming} computes a single Fisher information matrix $\mathbf{F}$ over all model parameters $\theta$ and penalizes deviations from the parameters $\theta^*$ learned on the previous task:
\begin{equation}
    \mathcal{L}_{\text{EWC}} = \sum_i F_i (\theta_i - \theta_i^*)^2.
    \label{eq:ewc_standard}
\end{equation}
This uniform treatment fails to account for the heterogeneous nature of multimodal encoders.
We observe empirically that structural encoder parameters exhibit high Fisher values and are most susceptible to forgetting due to topological drift between graph snapshots, while molecular encoder parameters---which process stable fingerprint features---show lower Fisher values and are more robust.

Modality-Aware EWC (MA-EWC) addresses this by partitioning model parameters into modality-specific groups $\theta = \{\theta_s, \theta_t, \theta_m, \theta_f\}$ (structural, textual, molecular, and fusion/decoder) and computing separate Fisher matrices for each:
\begin{equation}
    \mathcal{L}_{\text{MA-EWC}} = \lambda_s \sum_i F_i^s (\theta_{s,i} - \theta_{s,i}^*)^2 + \lambda_t \sum_i F_i^t (\theta_{t,i} - \theta_{t,i}^*)^2 + \lambda_m \sum_i F_i^m (\theta_{m,i} - \theta_{m,i}^*)^2 + \lambda_f \sum_i F_i^f (\theta_{f,i} - \theta_{f,i}^*)^2,
    \label{eq:maewc}
\end{equation}
where $\lambda_s{=}10$, $\lambda_t{=}5$, $\lambda_m{=}1$, and $\lambda_f{=}5$ are per-modality regularization strengths reflecting each encoder's forgetting susceptibility.
The Fisher matrices $\mathbf{F}^s, \mathbf{F}^t, \mathbf{F}^m, \mathbf{F}^f$ are computed after training on each task by accumulating squared gradients over the task's training data, following the diagonal Fisher approximation.

\subsection{Multimodal Memory Replay}
\label{subsec:replay}

To complement regularization-based forgetting prevention, \ours maintains a fixed-size memory buffer $\mathcal{M}$ of representative triples from previous tasks.
Na\"ive random selection risks overrepresenting dense graph regions while missing sparse but informative subgraphs.
We use K-means clustering on the structural embeddings of head entities to ensure diversity:
\begin{enumerate}
    \item After training on task $T_k$, extract structural embeddings $\{\mathbf{h}_v^s\}$ for all head entities in $T_k$'s training triples.
    \item Run K-means with $K = |\mathcal{M}| / k$ clusters (where $k$ is the task index) to partition the embedding space.
    \item Select the triple whose head entity is closest to each cluster centroid.
    \item Add selected triples to $\mathcal{M}$, evicting oldest entries if the buffer is full.
\end{enumerate}
During training on subsequent tasks, replay triples are sampled uniformly from $\mathcal{M}$ and mixed with the current task's training data.
The replay loss is:
\begin{equation}
    \mathcal{L}_{\text{replay}} = \frac{1}{|\mathcal{B}_r|} \sum_{(h,r,t) \in \mathcal{B}_r} \ell(f_r(\mathbf{h}_h, \mathbf{h}_t), y),
    \label{eq:replay}
\end{equation}
where $\mathcal{B}_r$ is a mini-batch of replayed triples and $\ell$ is the same loss function used for the current task.

\subsection{Training Procedure}
\label{subsec:training}

The complete \ours training procedure is summarized in \cref{alg:cmkl}.
For each task $T_k$ in the continual sequence, the model first encodes all entities through the three modality-specific encoders and fuses their representations via cross-modal attention.
The total training objective combines the task loss, MA-EWC regularization (for $k > 1$), and replay loss (for $k > 1$):
\begin{equation}
    \mathcal{L} = \mathcal{L}_{\text{task}} + \mathcal{L}_{\text{MA-EWC}} + \alpha \cdot \mathcal{L}_{\text{replay}},
    \label{eq:total}
\end{equation}
where $\alpha$ controls the replay weight.
After training on each task, per-modality Fisher matrices are computed and diverse exemplars are added to the memory buffer.

\begin{algorithm}[t]
\caption{\ours Training Procedure}
\label{alg:cmkl}
\begin{algorithmic}[1]
\REQUIRE Task sequence $\mathcal{T} = \{T_1, \ldots, T_K\}$, modality features $\mathbf{X}$, buffer size $|\mathcal{M}|$
\ENSURE Trained parameters $\theta = \{\theta_s, \theta_t, \theta_m, \theta_f\}$
\STATE Initialize encoders, fusion module, decoder
\STATE $\mathcal{M} \leftarrow \emptyset$ \COMMENT{Memory buffer}
\FOR{$k = 1$ \TO $K$}
    \STATE \textbf{// Encode entities with modality-specific encoders}
    \STATE $\mathbf{H}^s \leftarrow \text{R-GCN}(\mathcal{V}_k, \mathcal{E}_k; \theta_s)$
    \STATE $\mathbf{H}^t \leftarrow \text{BiomedBERT-Proj}(\mathbf{X}^{\text{text}}; \theta_t)$
    \STATE $\mathbf{H}^m \leftarrow \text{MLP}(\mathbf{X}^{\text{mol}}; \theta_m)$
    \STATE \textbf{// Fuse via cross-modal attention}
    \STATE $\mathbf{H} \leftarrow \text{CrossAttn}(\mathbf{H}^s, \mathbf{H}^t, \mathbf{H}^m; \theta_f)$
    \FOR{each epoch}
        \STATE $\mathcal{L}_{\text{task}} \leftarrow$ link prediction loss on $T_k$ using decoder($\mathbf{H}$)
        \IF{$k > 1$}
            \STATE $\mathcal{L}_{\text{MA-EWC}} \leftarrow$ \cref{eq:maewc} with stored Fisher matrices and $\theta^*$
            \STATE $\mathcal{L}_{\text{replay}} \leftarrow$ replay loss on sampled $\mathcal{B}_r \subset \mathcal{M}$
        \ELSE
            \STATE $\mathcal{L}_{\text{MA-EWC}} \leftarrow 0$; $\mathcal{L}_{\text{replay}} \leftarrow 0$
        \ENDIF
        \STATE $\mathcal{L} \leftarrow \mathcal{L}_{\text{task}} + \mathcal{L}_{\text{MA-EWC}} + \alpha \cdot \mathcal{L}_{\text{replay}}$
        \STATE Update $\theta$ via Adam optimizer
    \ENDFOR
    \STATE \textbf{// Post-task: compute per-modality Fisher and update buffer}
    \STATE Compute $\mathbf{F}^s, \mathbf{F}^t, \mathbf{F}^m, \mathbf{F}^f$ from squared gradients on $T_k$
    \STATE Store $\theta^* \leftarrow \theta$
    \STATE $\mathcal{M} \leftarrow \mathcal{M} \cup \text{K-means-Select}(\mathbf{H}^s, T_k, |\mathcal{M}|/k)$
\ENDFOR
\end{algorithmic}
\end{algorithm}
